% !PRINTED 685
\weberpage{685}
\S\ 196. \gesperrt{Geschlecht des K\"orpers $\Omega$.}

Da in einem Polygonquotienten Z\"ahler und Nenner von gleicher Ordnung sind, so folgt aus der letzten Formel [\S\ 195, (12)]:
\weq{1}{w_\alpha - 2a = w_\beta - 2b.}

Bezeichnen wir also mit $n$ die Ordnung und mit $w$ die Verzweigungszahl f\"ur eine beliebige Variable, die, wie wir im \S\ 193 gesehen haben, eine gerade Zahl ist, so ist die ganze Zahl:
\weq{2}{p = \tfrac{1}{2}w - n + 1,}
eine zu dem K\"orper $\Omega$ geh\"orige invariante ganze Zahl, die das \gesperrt{Geschlecht des K\"orpers} genannt wird. Da\ss{} diese Zahl nicht negativ sein kann, ergibt sich aus \S\ 193, (5) und (9), wonach
\weq{3}{p = (r_2-1)+(r_3-1)+\cdots+(r_{n_1}-1)}
% [sic] printed final subscript is $n_1$; formulas §193 (5), (9) require $n$.
ist, und so aus lauter Summanden besteht, deren keiner negativ ist.

Zwei Funktionen $\alpha$, $\beta$ nennen wir ein \gesperrt{primitives Paar} des K\"orpers $\Omega$, wenn alle Funktionen in $\Omega$ rational durch $\alpha$ und $\beta$ ausgedr\"uckt werden k\"onnen. Ist $\alpha$ von der $m$ten, $\beta$ von der $n$ten Ordnung, so ist nach \S\ 171 f\"ur ein primitives Paar notwendig und hinreichend, da\ss{} die zwischen $\alpha$ und $\beta$ bestehende irreduzible Gleichung:
\weq{4}{F(\overset{n}{\alpha},\overset{m}{\beta}) = 0}
in $\alpha$ vom $n$ten, in $\beta$ vom $m$ten Grade sei.

Es sollen die Funktionen $F'(\alpha)$, $F'(\beta)$ durch Polygonquotienten dargestellt werden. Wir setzen:
\weq{5}{\alpha = \dfrac{\mathfrak{M}}{\mathfrak{A}}, \quad \beta = \dfrac{\mathfrak{N}}{\mathfrak{B}},}
worin $\mathfrak{A}$, $\mathfrak{M}$ von der $m$ten, $\mathfrak{B}$, $\mathfrak{N}$ von der $n$ten Ordnung sind. $\mathfrak{A}$ hat dann mit $\mathfrak{M}$ und $\mathfrak{B}$ mit $\mathfrak{N}$ keinen Punkt gemein. Wir setzen:
\weq{6}{\begin{aligned}
F(\alpha,\beta) &= a_0\alpha^n + a_1\alpha^{n-1} + \cdots + a_n,\\
F'(\alpha) &= na_0\alpha^{n-1} + (n-1)a_1\alpha^{n-2} + \cdots + a_{n-1},\\
\alpha F'(\alpha) &= -a_1\alpha^{n-1} - 2a_2\alpha^{n-2} - \cdots - na_n.
\end{aligned}}

Aus der zweiten dieser Gleichungen schlie\ss{}t man, da\ss{} im Nenner von $F'(\alpha)$ keine Punkte vorkommen, die nicht in $\mathfrak{A}$ oder in $\mathfrak{B}$ enthalten sind, und aus dem dritten folgt, da\ss{} dieser Nenner h\"ochstens $= \mathfrak{A}^{n-2}\mathfrak{B}^m$ sein kann. Wir setzen also:
\weq{7}{F'(\alpha) = \dfrac{\mathfrak{L}}{\mathfrak{A}^{n-2}\mathfrak{B}^m}.}
